% Préambule commun aux cours IPR Scholar.
\documentclass[11pt,a4paper,oneside]{article}
\usepackage[T1]{fontenc}
\usepackage[utf8]{inputenc}
\usepackage[french]{babel}
\usepackage{lmodern}
\usepackage{microtype}
\usepackage[a4paper,margin=2.35cm]{geometry}
\usepackage{amsmath,amssymb,mathtools}
\usepackage{booktabs,longtable,array}
\usepackage{xcolor}
\usepackage{hyperref}
\usepackage{tikz}
\usetikzlibrary{arrows.meta,positioning,shapes.geometric,calc}
\hypersetup{colorlinks=true,linkcolor=blue!55!black,citecolor=blue!55!black,urlcolor=blue!55!black}
\definecolor{iprblue}{HTML}{1F4E79}
\definecolor{iprgreen}{HTML}{2D6A4F}
\newcommand{\IPR}{\textsc{ipr}}
\newcommand{\coursmeta}[2]{\author{Projet \IPR{} -- documentation Scholar}\date{#1}\maketitle\begin{center}\small #2\end{center}\vspace{0.5cm}\tableofcontents\clearpage}
\newcommand{\resultat}[1]{\paragraph{À retenir.} #1}

\title{Cours 1 -- Espace linéaire, luminance et gamut mapping}
\begin{document}
\coursmeta{3 octobre 2026}{Modèle utilisé avant, pendant et après tout calcul de rendu dans \IPR{}.}

\section{Problème traité}
Un fichier image n'est pas seulement une grille de triplets RGB : les valeurs peuvent être codées dans un espace colorimétrique non linéaire, avec ou sans profil ICC. Additionner ou filtrer directement ces valeurs encodées produit des résultats physiquement faux. \IPR{} convertit donc les entrées vers sRGB linéaire, travaille en flottants et n'encode en sRGB qu'à la sortie.

\section{Théorie}
La lumière est additive : doubler le nombre de photons double l'énergie. Une composante linéaire $r$ est donc proportionnelle à l'énergie, contrairement à une composante sRGB encodée qui privilégie la perception et l'efficacité de quantification. La luminance relative sRGB/Rec.~709 est
\begin{equation}
Y=0{,}2126r+0{,}7152g+0{,}0722b.
\end{equation}
Les coefficients sont normalisés et reflètent la plus grande sensibilité visuelle au vert~\cite{srgb,icc2022}. La conversion linéaire RGB--XYZ est une matrice $3\times3$ ; elle autorise ensuite les changements de primaires ou d'illuminant.

Un rendu peut quitter le cube sRGB. \IPR{} compresse alors la chromaticité autour du neutre de même luminance. Soit $\mathbf{c}=\mathbf{rgb}-Y(1,1,1)^T$. Le plus grand $s\in[0,1]$ tel que $Y\mathbf{1}+s\mathbf{c}\in[0,1]^3$ conserve la luminance et réduit uniquement l'excès de saturation.

\section{Application dans \IPR{}}
Les images standards sont d'abord converties par Qt en \texttt{Format\_RGBA64}, puis vers \texttt{SRgbLinear}; les valeurs sont normalisées par $65535$ dans un tampon \texttt{CV\_32FC4}. Les RAW 16~bits passent directement vers le même espace de travail. À l'export, la compression de gamut précède l'encodage sRGB et la quantification 8 ou 16~bits.

\section{Exemple}
Pour $\mathbf{p}=(0{,}50,0{,}25,0{,}10)$, $Y\simeq0{,}292$. Une exposition $+1$~EV doit donner environ $(1{,}00,0{,}50,0{,}20)$ en linéaire. Si l'on doublait les composantes sRGB encodées, le résultat ne correspondrait pas à un doublement d'énergie et les tons moyens seraient artificiellement trop lumineux.

\section{Limites et expérimentation}
La compression actuelle est une projection simple, pas un gamut mapping perceptuel CAM16. Elle est robuste et préserve $Y$, mais peut désaturer fortement une couleur HDR. Une évolution devra comparer cette stratégie à une compression dans OKLab ou JzAzBz sur un corpus à couleurs saturées.

\resultat{Les filtres, gains et flous doivent agir en linéaire ; l'espace sRGB est une destination d'affichage ou d'export, non un espace de calcul.}
\bibliographystyle{plain}\bibliography{references}
\end{document}
